\documentclass[aspectratio=169,10pt]{beamer}

\usefonttheme{professionalfonts}
\usepackage{fontspec}
\usepackage{unicode-math}
\usepackage{booktabs}
\usepackage{amsmath}
\usepackage{array}
\newcommand{\kw}[1]{\textcolor{navy}{\textbf{#1}}}

\setmainfont{texgyrepagella}[
  Extension      = .otf,
  UprightFont    = *-regular,
  BoldFont       = *-bold,
  ItalicFont     = *-italic,
  BoldItalicFont = *-bolditalic,
  Ligatures      = TeX,
  Scale          = 0.96
]
\setsansfont{texgyrepagella}[
  Extension      = .otf,
  UprightFont    = *-regular,
  BoldFont       = *-bold,
  ItalicFont     = *-italic,
  BoldItalicFont = *-bolditalic,
  Ligatures      = TeX,
  Scale          = 0.96
]
\setmathfont{texgyrepagella-math.otf}

\definecolor{ink}{HTML}{1C1C1C}
\definecolor{navy}{HTML}{1B3A4B}
\definecolor{rust}{HTML}{8C5A2E}
\definecolor{wash}{HTML}{F4F1EA}
\definecolor{muted}{HTML}{5C6670}
\definecolor{line}{HTML}{D9D3C7}

\setbeamercolor{normal text}{fg=ink,bg=white}
\setbeamercolor{structure}{fg=navy}
\setbeamercolor{frametitle}{fg=navy,bg=white}
\setbeamercolor{title}{fg=navy}
\setbeamercolor{subtitle}{fg=navy}
\setbeamercolor{author}{fg=ink}
\setbeamercolor{date}{fg=muted}
\setbeamercolor{itemize item}{fg=rust}
\setbeamercolor{itemize subitem}{fg=rust}
\setbeamercolor{enumerate item}{fg=rust}
\setbeamercolor{footline}{fg=muted,bg=white}
\setbeamercolor{block title}{fg=white,bg=navy}
\setbeamercolor{block body}{fg=ink,bg=wash}

\setbeamerfont{frametitle}{series=\bfseries,size=\large}
\setbeamerfont{footline}{size=\scriptsize}
\setbeamertemplate{navigation symbols}{}
\setbeamertemplate{itemize item}{\small\color{rust}$\bullet$}
\setbeamertemplate{itemize subitem}{\scriptsize\color{rust}$\bullet$}
\setbeamersize{text margin left=8mm,text margin right=8mm}

\setbeamertemplate{frametitle}{%
  \vspace{0.15em}%
  {\usebeamerfont{frametitle}\color{navy}\insertframetitle}\par
  \vspace{0.15em}%
  {\color{rust}\rule{\textwidth}{0.7pt}}\par
  \vspace{0.15em}%
}
\newcommand{\currentsource}{}
\newcommand{\framesource}[1]{\gdef\currentsource{#1}}
\newcommand{\aiassisted}{AI assisted, via Grok 4.7}

\setbeamertemplate{footline}{%
  \leavevmode
  \hbox to \paperwidth{%
    \hfil
    \begin{beamercolorbox}[wd=\dimexpr\paperwidth-16mm\relax,ht=2.5ex,dp=1.1ex,right]{footline}%
      \currentsource
    \end{beamercolorbox}%
    \hspace{8mm}%
  }%
}

\begin{document}

\begin{frame}[plain]
\vspace{1.1em}
{\small\scshape\color{rust}APMA 4903 \quad 4 November 2026}\par\vspace{0.45em}
{\LARGE\bfseries\color{navy}The cost of a forced sale}\par\vspace{0.35em}
{\large\color{navy}When the book has not refilled before the next trade.}\par
\vspace{0.9em}
{\color{rust}\rule{3.2cm}{1.1pt}}\par\vspace{0.7em}
Daniel Holland\par
{\small\color{muted}with Allison Duffy and Jasmine Li}\par
\vspace{1.1em}
{\small The reserve needs the expected bill for selling a dollar shortfall.\\
This section prices that sale, and checks the schedule with Q-learning.}
\end{frame}

\begin{frame}{Where this section sits}
\framesource{Allison Duffy and Jasmine Li; Almgren and Chriss, 2001}
\begin{itemize}
\item The shortfall is \(Q=\max\{L-R,0\}\) dollars. Selling it means liquidating \(X\approx Q/S_0\) shares before the call is due.
\item Allison prices that liquidation with Almgren--Chriss, at risk aversion zero, and minimizes cash carry plus expected sale cost.
\item Jasmine turns a sale cost per dollar, \(\kappa\), into a tolerance for how often the reserve may be short.
\item The bill below is the input those two formulas use. The half-life of the book is the part their refill formula leaves out.
\end{itemize}
\end{frame}

\begin{frame}{Three learning problems}
\framesource{Sutton and Barto, 2018, \S 1.3}
\begin{columns}[T]
\begin{column}{0.32\textwidth}
\textbf{\color{navy}Supervised}\\[0.35em]
Examples come in pairs: an input, and the correct output.\\[0.4em]
The error is the distance to that output.
\end{column}
\begin{column}{0.32\textwidth}
\textbf{\color{navy}Unsupervised}\\[0.35em]
Inputs arrive with no target output.\\[0.4em]
The job is to find structure in the inputs themselves.
\end{column}
\begin{column}{0.32\textwidth}
\textbf{\color{navy}Reinforcement}\\[0.35em]
No correct action is supplied.\\[0.4em]
The agent acts, and the environment answers with a single number.
\end{column}
\end{columns}
\vspace{0.9em}
The number scores the step. The object is the total of those numbers over the horizon, and the record of what happened depends on the actions already taken.
\end{frame}

\begin{frame}{What makes the third one different}
\framesource{Sutton and Barto, 2018, \S 1.3}
\begin{itemize}
\item Supervised learning can be told, after the fact, which trade would have been best. This sale cannot. The fund sees the bill of the trade it chose.
\item The same state can be followed by a different bill next time, because the book and the price move.
\item A cheap step can be a bad decision. Selling nothing today has no walk cost, and it leaves the whole position for the deadline.
\item The data are produced by the policy. A different schedule would have left a different displacement, and the next sample would have been a different one.
\end{itemize}
\end{frame}

\begin{frame}{A policy}
\framesource{Sutton and Barto, 2018, \S 1.3. Wording is not a quotation.}
\begin{itemize}
\item A policy maps the state the agent sees to the action it takes.
\item It may be a formula, a table, or the result of a search. It may be random. The policy alone decides what the agent does.
\item In this sale the state on a bar is the inventory still held and the displacement of the book. The action is how many lots to sell.
\item The table \(n_k(x,y)\) is the policy. Once it is filled, trading does not require another optimization.
\end{itemize}
\end{frame}

\begin{frame}{Reward, and value}
\framesource{Sutton and Barto, 2018, \S 1.3. Wording is not a quotation.}
\begin{itemize}
\item The reward is the number attached to the current step. The environment computes it from the state and the action. The agent changes it only by acting.
\item This problem is written as a cost. The stage cost is the negative of a reward, and the update takes a minimum where Sutton and Barto take a maximum.
\item The value of a state is the expected remaining total, from that state to the deadline.
\item A step can be cheap and the state still poor, because of the states that follow. Decisions use values. The costs themselves are observed; the values are estimated from the sequence of sales.
\end{itemize}
\end{frame}

\begin{frame}{A model, and doing without one}
\framesource{Sutton and Barto, 2018, \S 1.3. Wording is not a quotation.}
\begin{itemize}
\item A model predicts the next state and the next cost from the current state and the action.
\item Planning chooses the action by using that prediction before the step is experienced. Methods that do this are model-based.
\item A model-free method never receives the prediction rule. It receives the cost that occurred and the state that came next.
\item The dynamic program plans with the decay rule. Q-learning uses only the sampled cost and the sampled next displacement.
\end{itemize}
\end{frame}

\begin{frame}{The four objects in this sale}
\framesource{Sutton and Barto, 2018, \S 1.3. Sale column: \aiassisted.}
\small
\begin{tabular}{@{}lp{0.72\textwidth}@{}}
\toprule
Object & In the forced sale \\
\midrule
State & Day, shares still held, displacement of our own selling \\
Action & Lots sold on this bar \\
Stage cost & Cash the walk gives up, plus a price-risk penalty if risk aversion is positive \\
Value & Expected remaining cost through the deadline \\
Model & Depth \(q\) and half-life \(h\), through \(\beta=2^{-\tau/h}\) \\
Policy & Lots to sell in each state. The greedy table is enough to trade \\
\bottomrule
\end{tabular}
\vspace{0.7em}

The reserve uses risk aversion zero, so the stage cost on the slides below is the walk alone. The price-risk penalty is Almgren and Chriss's other objective.
\end{frame}

\begin{frame}{The sale}
\framesource{Almgren and Chriss, 2001}
Sell \(X\) shares over \(N\) days, \(T=N\tau\), arrival price \(S_0\).
\[
x_0=X,\qquad x_N=0,\qquad n_k=x_{k-1}-x_k\geq 0.
\]
The unaffected price is the Almgren--Chriss martingale, without their permanent-impact term,
\[
S_k=S_{k-1}+\sigma\sqrt{\tau}\,\xi_k.
\]
\begin{itemize}
\item The dollar shortfall converts at the arrival price, \(X=Q/S_0\), as on the reserve slide.
\item A fee \(\varepsilon X\) and a permanent term \(\gamma X^{2}/2\) are added to the bill. They are the same for every completed schedule, so they do not change the trades.
\item The tables set both to zero. The slides leave them symbolic.
\end{itemize}
\end{frame}

\begin{frame}{A book that remembers}
\framesource{Obizhaeva and Wang, 2013, Proposition 3}
Depth \(q\) shares per dollar. Selling \(n\) shares, starting from a displacement \(y\), costs
\[
n\Bigl(y+\frac{n}{2q}\Bigr).
\]
A fraction of that displacement is still there next day,
\[
\beta=2^{-\tau/h}.
\]
\begin{itemize}
\item \(h\) is the half-life. It is a property of the public book.
\item Half a day left means \(\beta=0.5\): half of today's sale is still in the book tomorrow.
\item The fund is assumed to see \(y\) before it trades.
\end{itemize}
\end{frame}

\begin{frame}{Refill is a special case}
\framesource{Almgren and Chriss, 2001, eqs.~(9)--(10)}
\begin{itemize}
\item Almgren--Chriss, at risk aversion zero, sells the same number of shares each day.
\item Their temporary cost matches the walk only when the book is flat at the start of every bar. That is \(\beta=0\): the book refills before the next trade.
\item Matching the two formulas gives their temporary coefficient, \(\eta=\tau/(2q)\).
\item The risk-neutral refill cost is then
\[
E_{\min}(X)=\varepsilon X+\Bigl(\frac{\gamma}{2}+\frac{\eta}{T}\Bigr)X^{2}.
\]
\item With \(\varepsilon=\gamma=0\) and ten days, that bill is \(X^{2}/(2qN)\). On the illustration it is \(5.00\) basis points.
\end{itemize}
\end{frame}

\begin{frame}{Implementation shortfall}
\framesource{Perold, 1988; Almgren and Chriss, 2001}
Perold's shortfall is arrival value minus cash received.
\begin{itemize}
\item If the trade list ignores the unaffected price, the expected shortfall equals the expected walk.
\item The price-risk standard deviation is \(\sigma\sqrt{\tau\sum_k x_k^{2}}\), divided by arrival value.
\item At risk aversion zero this standard deviation does not enter the bill. It is large, because a multi-day sale still holds a large position.
\item On the one-day book the best schedule accepts a bit more price risk than equal slices, in exchange for a smaller walk.
\end{itemize}
\end{frame}

\begin{frame}{The bill the reserve uses}
\framesource{Almgren and Chriss, 2001; Allison Duffy and Jasmine Li}
\[
C(Q)
=\varepsilon X+\frac{\gamma}{2}X^{2}
+\mathbb{E}\!\left[\sum_{k=1}^{N}n_k\Bigl(y_{k-1}+\frac{n_k}{2q}\Bigr)\right],
\qquad
\kappa(Q)=\frac{C(Q)}{Q}.
\]
\begin{itemize}
\item \(\kappa\) is dollars of sale cost per dollar of shortfall. The tables quote it in basis points.
\item Allison minimizes \(cR+\mathbb{E}[C(Q(R))]\), with \(c\) the cost of holding a dollar of cash.
\item Jasmine's flat-cost rule is the special case \(C(q)=\kappa q\), which gives \(\varepsilon^{\star}=\rho/\kappa\). Her \(\rho\) is Allison's \(c\).
\end{itemize}
\end{frame}

\begin{frame}{A quadratic bill has a different derivative}
\framesource{\aiassisted}
For \(C(q)=aq+bq^{2}\),
\[
\frac{d}{dR}\,\mathbb{E}[C(Q(R))]
=-a\,\Pr(L>R)-2b\,\mathbb{E}[Q].
\]
\begin{itemize}
\item The flat-cost tolerance uses only the probability of a shortfall, and only when \(b=0\).
\item The walk, and her permanent-impact term, are quadratic. The derivative then uses the probability and the expected shortfall.
\item A single \(\kappa\) is still the right input to \(\varepsilon^{\star}=\rho/\kappa\), read at a stated shortfall. The curve \(C(Q)\) is the input to Allison's objective.
\end{itemize}
\end{frame}

\begin{frame}{The dynamic program}
\framesource{\aiassisted}
\(V_k(x,y)\) is the remaining cost with inventory \(x\) and displacement \(y\). The last trade liquidates:
\[
V_{N-1}(x,y)=x\Bigl(y+\frac{x}{2q}\Bigr).
\]
Earlier,
\[
V_k(x,y)
=\min_{0\leq n\leq x}
\left\{
n\Bigl(y+\frac{n}{2q}\Bigr)
+\lambda\sigma^{2}\tau(x-n)^{2}
+\mathbb{E}[V_{k+1}(x-n,y')]
\right\}.
\]
\begin{itemize}
\item The expectation uses the decay rule. This is the model-based calculation.
\item The reserve sets \(\lambda=0\). The variance term is then absent, and the schedule minimizes the expected walk.
\end{itemize}
\end{frame}

\newsavebox{\dpalg}
\sbox{\dpalg}{%
\footnotesize
\renewcommand{\arraystretch}{1.18}%
\begin{tabular}{@{}r@{\hspace{0.55em}}l@{}}
 & \kw{Input:} depth \(q\), decay \(\beta\), risk aversion \(\lambda\) \\
{\color{muted}1} & \kw{for} every inventory \(x\) and displacement \(y\) \kw{do} \\
{\color{muted}2} & \quad \(V_{N-1}(x,y) \gets x\bigl(y + x/(2q)\bigr)\) \\
{\color{muted}3} & \quad \(n_{N-1}(x,y) \gets x\) \\
{\color{muted}4} & \kw{for} \(k = N-2\) down to \(0\) \kw{do} \\
{\color{muted}5} & \quad \kw{for} every inventory \(x\) and displacement \(y\) \kw{do} \\
{\color{muted}6} & \quad\quad \kw{for} every legal sale \(n = 0,\ldots,x\) \kw{do} \\
{\color{muted}7} & \quad\quad\quad \(c \gets n\bigl(y + n/(2q)\bigr) + \lambda\sigma^{2}\tau(x-n)^{2}\) \\
{\color{muted}8} & \quad\quad\quad \(y' \gets \beta\bigl(y + n/q\bigr)\), plus a refill shock if the book is noisy \\
{\color{muted}9} & \quad\quad\quad score the sale by \(c + \mathbb{E}\bigl[V_{k+1}(x-n, y')\bigr]\) \\
{\color{muted}10} & \quad\quad \(V_k(x,y) \gets\) the smallest score, and \(n_k(x,y)\) the sale that attains it \\
\end{tabular}}

\newsavebox{\qalg}
\sbox{\qalg}{%
\footnotesize
\renewcommand{\arraystretch}{1.22}%
\begin{tabular}{@{}r@{\hspace{0.55em}}l@{}}
 & \kw{Input:} sampled costs and next displacements, not \(\beta\) and not \(q\) \\
{\color{muted}1} & set every action value to \(+\infty\) \\
{\color{muted}2} & \kw{for} every inventory \(x\) and displacement \(y\) \kw{do} \\
{\color{muted}3} & \quad sell what remains, and store the observed cost as \(Q_{N-1}(x,y,x)\) \\
{\color{muted}4} & \kw{for} \(k = N-2\) down to \(0\) \kw{do} \\
{\color{muted}5} & \quad \kw{for} every \(x\), \(y\), and legal sale \(n\) \kw{do} \\
{\color{muted}6} & \quad\quad take the sale; observe the cost \(c\) and the next displacement \(y'\) \\
{\color{muted}7} & \quad\quad \(Q_k(x,y,n) \gets c + \min_{n' \le x-n} Q_{k+1}(x-n, y', n')\) \\
{\color{muted}8} & \(n_k(x,y) \gets \arg\min_{n} Q_k(x,y,n)\) \\
\end{tabular}}

\newsavebox{\tradealg}
\sbox{\tradealg}{%
\footnotesize
\renewcommand{\arraystretch}{1.22}%
\begin{tabular}{@{}r@{\hspace{0.55em}}l@{}}
{\color{muted}1} & start from the full shortfall, \(x \gets X\), and a flat book, \(y \gets 0\) \\
{\color{muted}2} & \kw{for} \(k = 0\) to \(N-1\) \kw{do} \\
{\color{muted}3} & \quad read today's displacement \(y\) \\
{\color{muted}4} & \quad sell \(n \gets n_k(x,y)\) lots \\
{\color{muted}5} & \quad on the last day, sell whatever inventory remains \\
{\color{muted}6} & \quad \(x \gets x - n\) \\
{\color{muted}7} & \quad take the displacement the book returns as tomorrow's \(y\) \\
\end{tabular}}

\begin{frame}{The dynamic program, as a procedure}
\framesource{\aiassisted}
\usebox{\dpalg}

\vspace{0.7em}
Line 8 is the model: the next displacement is computed from the decay rule. The reserve sets \(\lambda=0\), so line 7 is the walk alone. Line 9 then adds the expected cost of whatever inventory is left.
\end{frame}

\begin{frame}{Two limits the grid has to match}
\framesource{Almgren and Chriss, 2001; Bertsimas and Lo, 1998; Obizhaeva and Wang, 2013}
\begin{itemize}
\item \(\beta=0\) and \(\lambda=0\): equal trades, and the bill \(X^{2}/(2qN)\). This is Almgren--Chriss, and Bertsimas and Lo for a quadratic temporary cost.
\item \(\beta=0\) and \(\lambda>0\): the hyperbolic holdings. On the illustration, the reference penalty raises the refill bill from \(5.00\) to \(8.12\) basis points. The reserve does not use that penalty.
\item A finite half-life and \(\lambda=0\): more shares on the first day and the last day than on the days between. That shape is Obizhaeva and Wang.
\item If the book never refills, every completed schedule costs \(X^{2}/(2q)\).
\end{itemize}
\end{frame}

\begin{frame}{The state the learner sees}
\framesource{\aiassisted}
\begin{itemize}
\item Inventory is an integer number of lots. Displacement is rounded to a grid of nodes, from a flat book up to a one-shot walk.
\item The dynamic program and the learner solve that same finite Markov decision process.
\item The action value \(Q_k(x,y,n)\) is the cost of selling \(n\) lots now and acting optimally afterward.
\item In this section \(Q_k\) is that action value. In the reserve, \(Q\) is the dollar shortfall. They do not appear in one equation.
\end{itemize}
\end{frame}

\begin{frame}{Q-learning, filled from the deadline}
\framesource{Watkins and Dayan, 1992. Finite-horizon sweep: \aiassisted.}
Watkins's update replaces the expectation by a sample, with step size \(1/(\text{number of visits})\):
\[
Q_k(x,y,n)
\leftarrow
Q_k(x,y,n)
+\alpha\Bigl(
c+\min_{n'}Q_{k+1}(x',y',n')
-Q_k(x,y,n)
\Bigr).
\]
\begin{itemize}
\item The last bar has no continuation. Its value is the observed cost of selling what remains.
\item Each earlier bar is sampled only after the next bar has been filled, so the target has stopped moving.
\item With no refill shock, one visit to each state and action is the conditional expectation. The sweep reproduces the dynamic program.
\item The learner is not given \(\beta\) or \(q\). It is given the sampled cost and the sampled next displacement.
\end{itemize}
\end{frame}

\begin{frame}{Q-learning, as a procedure}
\framesource{Watkins and Dayan, 1992. Procedure: \aiassisted.}
\usebox{\qalg}

\vspace{0.7em}
Line 6 replaces the decay rule. Several samples of one state and action are averaged, which is Watkins's update with step size \(1/(\text{number of visits})\). With no refill shock, one sample is the conditional expectation, and the sale on line 8 is the dynamic program's sale.
\end{frame}

\begin{frame}{From the table to a trade}
\framesource{\aiassisted}
\usebox{\tradealg}

\vspace{0.7em}
\begin{itemize}
\item The same loop plays either table. Only \(n_k(x,y)\) changes.
\item The dynamic program filled its table from \(\beta\) and \(q\). The learner filled its table from the samples.
\end{itemize}
\end{frame}

\begin{frame}{Why the horizon is filled backward}
\framesource{Watkins and Dayan, 1992. Finite-horizon argument: \aiassisted.}
\begin{itemize}
\item The same update, run forward from an unfinished table, treats an action that has never been tried as free, because the minimum prefers a zero that has not been sampled.
\item That bias stays until every legal successor has been visited. On this grid it does not clear in a practical number of episodes.
\item Backward order removes it. A bar is sampled only after every legal action on the next bar has a value.
\item Watkins and Dayan prove almost-sure convergence for a discounted infinite horizon. The claim used here is the finite-horizon argument above. The dollars are not discounted.
\end{itemize}
\end{frame}

\begin{frame}{Why this algorithm}
\framesource{Mnih et al., 2015; Schulman et al., 2017; Haarnoja et al., 2018}
\begin{itemize}
\item The fixed point of the update is the minimizer of the Bellman equation. That is the claim we need: the schedule is optimal for this process when the decay rule is sampled rather than supplied.
\item A deep Q-network stores the action value in a neural net and trains it with a discounted target. The net is what makes the method run when the state will not fit in a table, as in Ning, Lin, and Jaimungal. Here the state is two integers. A net would replace the equality with the dynamic program by an approximation error.
\item Proximal policy optimization updates a parameterized policy from sampled returns. It does not build the action value, so there is no cell to compare with the dynamic program.
\item Soft actor-critic adds the entropy of a continuous action. That is a different objective, and the action here is a whole number of lots.
\end{itemize}
\end{frame}

\begin{frame}{What is gained past Almgren--Chriss}
\framesource{\aiassisted}
\begin{itemize}
\item Almgren--Chriss is the optimum under one model: impact is gone before the next trade, and the temporary coefficient \(\eta\) is known.
\item If that model is wrong, the formula still quotes \(5.00\) basis points. On a one-day half-life the sale costs \(10.81\).
\item The dynamic program repairs the model, provided the half-life and the depth are known and correct.
\item Q-learning repairs the schedule when those two numbers are not handed over. The observed walk and the observed displacement replace them.
\item In these runs the observations are drawn from the book the dynamic program uses, so the two can be compared. They agree to the cent. That agreement is the evidence. It is not a second bill.
\end{itemize}
\end{frame}

\begin{frame}{Hendricks and Wilcox, for scale}
\framesource{Hendricks and Wilcox, 2014}
\begin{itemize}
\item They also fill a table with Watkins's update.
\item Their action multiplies a fixed Almgren--Chriss slice by a factor between \(0\) and \(2\). The book is assumed to have recovered before the next bar.
\item The action here is the number of lots. The state their table does not carry is the displacement those lots leave behind.
\item The extension is that state. The learner is then in a position to sell less into a book it has already pushed down.
\end{itemize}
\end{frame}

\begin{frame}{Letters that are shared with the other sections}
\framesource{Allison Duffy and Jasmine Li; Almgren and Chriss, 2001}
\begin{itemize}
\item Almgren--Chriss \(\lambda\) is risk aversion, in units of one per dollar. The reserve sets it to zero.
\item Jasmine's \(\lambda\) is the Poisson rate at which a buyer appears for the private stake.
\item Jasmine's \(\rho\), and Allison's \(c\), are the cost of holding cash.
\item The speed of the public book is the half-life \(h\), reported as \(\beta=2^{-\tau/h}\).
\item Her crisis slows the Poisson rate. The downturn below changes the public book the shares are sold into.
\end{itemize}
\end{frame}

\begin{frame}{The illustration}
\framesource{\aiassisted}
\begin{itemize}
\item \(Q=5\) million dollars, sold as \(X=100{,}000\) shares at \(S_0=50\). One basis point is \(500\) dollars.
\item Ten trading days. Normal depth \(q=200{,}000\) shares per dollar. Daily volatility one dollar.
\item Selling the whole block at once walks the mid by \(50\) cents. The average fill is \(25\) cents worse: \(50\) basis points.
\item Risk aversion, the fee, and permanent impact are zero. Then \(C\) is the expected walk and \(\kappa=C/Q\).
\item Grid: \(40\) lots and \(41\) displacement nodes. These rows are \texttt{execution/report.py}.
\end{itemize}
\end{frame}

\newsavebox{\kappatable}
\sbox{\kappatable}{%
\footnotesize
\begin{tabular}{@{}lrrrr@{}}
\toprule
& \multicolumn{4}{c}{\(\kappa\), basis points} \\
\cmidrule(lr){2-5}
Half-life & Her formula & Equal trades & Optimal & Learner \\
\midrule
within the day & 5.00 & 5.00 & 5.00 & 5.00 \\
0.5 day & 5.00 & 7.25 & 7.25 & 7.25 \\
1 day & 5.00 & 11.50 & 10.81 & 10.81 \\
2 days & 5.00 & 20.25 & 17.94 & 17.94 \\
5 days & 5.00 & 32.75 & 27.50 & 27.50 \\
10 days & 5.00 & 39.50 & 34.25 & 34.25 \\
\bottomrule
\end{tabular}%
}

\newsavebox{\downturntable}
\sbox{\downturntable}{%
\footnotesize
\begin{tabular}{@{}lrrrr@{}}
\toprule
& \multicolumn{4}{c}{\(\kappa\), basis points} \\
\cmidrule(lr){2-5}
Downturn book & Her formula & Equal trades & Optimal & Learner \\
\midrule
refill assumed & 10.00 & 10.00 & 10.00 & 10.00 \\
half-life 5 days & 10.00 & 65.50 & 55.00 & 55.00 \\
\bottomrule
\end{tabular}%
}

\newsavebox{\sizetable}
\sbox{\sizetable}{%
\footnotesize
\begin{tabular}{@{}lrrr@{}}
\toprule
& \multicolumn{3}{c}{Optimal \(\kappa\), basis points} \\
\cmidrule(lr){2-4}
Shortfall & Her formula, 1 day & Optimal, 1 day & Optimal, downturn \\
\midrule
1.25 million & 1.25 & 2.70 & 13.75 \\
2.5 million & 2.50 & 5.41 & 27.50 \\
5 million & 5.00 & 10.81 & 55.00 \\
10 million & 10.00 & 21.62 & 110.00 \\
\bottomrule
\end{tabular}%
}

\begin{frame}{Cost per dollar of shortfall}
\framesource{\aiassisted}
\usebox{\kappatable}

\vspace{0.7em}
Her formula assumes the book refills, so it does not depend on the half-life. The learner repeats the optimal column. At a half-day the grid does not improve on equal trades.
\end{frame}

\begin{frame}{One day of memory}
\framesource{\aiassisted}
\begin{itemize}
\item Fraction left overnight: \(\beta=0.50\).
\item Refill formula: \(5.00\) basis points, \(2{,}500\) dollars.
\item Equal slices in the book that actually persists: \(11.50\) basis points, \(5{,}750\) dollars.
\item Best schedule: \(10.81\) basis points, \(5{,}406\) dollars. Per dollar of shortfall, \(\kappa=0.001081\).
\item The larger gap is persistence, \(5.00\) against \(10.81\). The smaller gap is reoptimization, \(11.50\) against \(10.81\).
\item Price-risk standard deviation: \(337.64\) basis points for equal slices, \(356.72\) for the best schedule.
\end{itemize}
\end{frame}

\begin{frame}{The sweep on that book}
\framesource{\aiassisted}
\begin{itemize}
\item Before any earlier bar is sampled, the greedy rule sells nothing until the deadline. The bill is the one-shot walk, \(25{,}000\) dollars.
\item Each completed bar lowers that cost. With no refill shock, the last point equals the dynamic program.
\item On the fine grid the curve ends at \(5{,}406\) dollars. The gap is zero.
\item The horizontal line on the screen is that one number. It is the optimum of the persistent book, not the \(2{,}500\)-dollar refill formula.
\item The column headed Learner is this same check, on every half-life.
\end{itemize}
\end{frame}

\begin{frame}{The public book in a downturn}
\framesource{\aiassisted}
Same shortfall, same ten days, same risk aversion. Depth halved, half-life five days, volatility two dollars.

\vspace{0.45em}
\usebox{\downturntable}

\vspace{0.55em}
\begin{itemize}
\item At risk aversion zero, volatility does not enter \(C\). Halving the depth doubles the refill formula, from \(5\) to \(10\).
\item The half-life takes the best schedule to \(55.00\) basis points, \(27{,}500\) dollars, \(\kappa=0.0055\).
\item The same half-life at the original depth costs \(27.50\). The ratio is the ratio of the depths.
\end{itemize}
\end{frame}

\begin{frame}{The cost per dollar grows with the shortfall}
\framesource{\aiassisted}
\usebox{\sizetable}

\vspace{0.7em}
Doubling the shortfall doubles \(\kappa\). The walk is quadratic, so a constant \(\kappa\) belongs to one shortfall. Allison's expectation uses the column.
\end{frame}

\begin{frame}{What is handed across}
\framesource{\aiassisted}
\begin{itemize}
\item To the reserve formula that assumes refill: replace \((\eta/T)X^{2}\) by the expected walk of the best schedule. Add \(\varepsilon X+(\gamma/2)X^{2}\) when those parameters are chosen.
\item To the flat-cost tolerance: \(\kappa=C(Q)/Q\) at a stated \(Q\) and a stated book. At five million dollars, \(0.001081\) in the normal one-day book and \(0.0055\) in the downturn book.
\item To the day the reserve is short: sell by the feedback rule \(n_k(x,y)\), not by the equal-trade list.
\item Q-learning is the check that this rule can be obtained from observed sales. The number in the reserve is \(C\), and on these runs the learner and the dynamic program return the same \(C\).
\end{itemize}
\end{frame}

\begin{frame}{References}
\framesource{}
\footnotesize
\begin{enumerate}
\item R.~Almgren and N.~Chriss, Optimal execution of portfolio transactions, \emph{Journal of Risk} 3(2):5--39, 2001. doi:10.21314/JOR.2001.041. Reserve formula: their eqs.~(9)--(10), p.~12.
\item D.~Bertsimas and A.~W. Lo, Optimal control of execution costs, \emph{Journal of Financial Markets} 1(1):1--50, 1998.
\item A.~A. Obizhaeva and J.~Wang, Optimal trading strategy and supply/demand dynamics, \emph{Journal of Financial Markets} 16(1):1--32, 2013. Proposition~3.
\item A.~F. Perold, The implementation shortfall: paper versus reality, \emph{Journal of Portfolio Management} 14(3):4--9, 1988.
\item R.~S. Sutton and A.~G. Barto, \emph{Reinforcement Learning: An Introduction}, 2nd ed., MIT Press, 2018, \S 1.3.
\item C.~J.~C.~H. Watkins and P.~Dayan, Q-learning, \emph{Machine Learning} 8:279--292, 1992.
\end{enumerate}
\end{frame}

\begin{frame}{References}
\framesource{}
\footnotesize
\begin{enumerate}
\setcounter{enumi}{6}
\item D.~Hendricks and D.~Wilcox, A reinforcement learning extension of the Almgren--Chriss framework, IEEE CIFEr, 2014. The PDF is in this repository.
\item V.~Mnih, K.~Kavukcuoglu, D.~Silver, et al., Human-level control through deep reinforcement learning, \emph{Nature} 518:529--533, 2015.
\item B.~Ning, F.~H.~Y. Lin, and S.~Jaimungal, Double deep Q-learning for optimal execution, arXiv:1812.06600, 2018. \emph{Applied Mathematical Finance}, 2021.
\item J.~Schulman, F.~Wolski, P.~Dhariwal, A.~Radford, and O.~Klimov, Proximal policy optimization algorithms, arXiv:1707.06347, 2017.
\item T.~Haarnoja, A.~Zhou, P.~Abbeel, and S.~Levine, Soft actor-critic, ICML, PMLR 80:1861--1870, 2018.
\end{enumerate}
\vspace{0.8em}
{\small The math note is \texttt{notes/execution-math.tex}. The screen is \texttt{streamlit run execution/app.py}.}
\end{frame}

\end{document}
